%------------------------------------------------------------------------------%
\section{Convergência da Série de Taylor}

Uma consequência da definição da Série de Taylor é que se uma função pode ser
representada por uma Série de Potências essa série coincide com
a Série de Taylor da função. Esse fato é apresentado com precisão pela
proposição a seguir.

%------------------------------------------------------------------------------%
\begin{proposicao}{Igualdade das Séries de Potências e Taylor}{prop:serietaylorpotencia}
  Se a função $f$ for igual a uma Série de Potências,
  no intervalo de convergência da série,
  \[
    f(x)=\sum_{n=0}^\infty c_n(x-a)^n
  \]
  então os coeficientes da série são os coeficientes de Taylor
  \[
    c_n = \frac{f^{(n)}(a)}{n!} \qquad n = 0, 1, 2, \dots
  \]
\end{proposicao}
%------------------------------------------------------------------------------%
\begin{proof}
  A função $f$ é infinitamente derivável, pois as séries de potências
  são infinitamente deriváveis em seu intervalo de convergência.
  Calculando a $n$-ésima derivada de $f$, temos
  \begin{align*}
    f^{(n)}(x) & = \sum_{k=n}^\infty k(k-1)\cdots(k-n+1)c_k(x-a)^{k-n}            \\
               & = n!c_n + \sum_{k=n+1}^\infty k(k-1)\cdots(k-n+1)c_k(x-a)^{k-n}
  \end{align*}
  Avaliando a derivada em $a$ temos
  \[
    f^{(n)}(a) = n!c_n
  \]
  Isolando $c_n$ concluímos a demonstração
  \[
    c_n = \frac{\;f^{(n)}(a)\;}{n!}
  \]
\end{proof}
%------------------------------------------------------------------------------%

A utilidade dessa última proposição é limitada, pois precisamos saber de antemão
que a função pode ser escrita como uma Série de Potências e essa verificação não
é trivial. Por exemplo, existem funções infinitamente deriváveis,
cujas séries de Taylor são convergentes mas não convergem para a função.
O próximo exemplo apresenta uma dessas funções.

%------------------------------------------------------------------------------%
\begin{example}
  Calcule a Série de Taylor centrada na origem, $a=0$, da função
  \[
    f(x) = \begin{cases}
      0                 & x \leq 0 \\[2mm]
      e^{\sfrac{-1}{x}} & x > 0
    \end{cases}
  \]
  Calcule a soma da série e compare com a função.
  \solution
  Essa é uma função suave, infinitamente derivável, cujo gráfico está
  ilustrado na figura a seguir.
  %---------------------------------------------------------------------------%
  \IncludeTexGraphicCentred{fig_functions_exp_1dx}[-0.6][-2]
  %---------------------------------------------------------------------------%
  Primeiro calculamos as derivadas de $f$ em zero. Como a função é definida
  com expressões diferentes de cada lado da origem precisamos calcular as derivadas
  de cada lado e verificar que ambas coincidem em $x=0$.

  Calcular as derivadas do lado esquerdo, $x<0$, é trivial pois a função é constante
  e portanto todas as suas derivadas são nulas. Consequentemente, o limite
  na origem também é zero
  \[
    \lim_{x\to 0^-} f^{(k)}(x) = 0 \qquad k = 1, 2, 3, \dots
  \]
  Do lado direito, $x>0$, a função tem a forma \(f(x) = e^{\sfrac{-1}{x}}\), calculando as
  primeiras derivadas obtemos
  \begin{align*}
    f'(x)      & = \frac{e^{\sfrac{-1}{x}}}{x^2}                       \\[3mm]
    f''(x)     & = \frac{e^{\sfrac{-1}{x}}}{x^4}\,(1-2x)               \\[3mm]
    f^{(3)}(x) & = \frac{e^{\sfrac{-1}{x}}}{x^6}\,(6x^2-6x+1)          \\[3mm]
    f^{(4)}(x) & = \frac{e^{\sfrac{-1}{x}}}{x^8}\,(-24x^3+36x^2-12x+1)
  \end{align*}
  Observamos, sem demonstrar, que a $n$-ésima derivada será sempre a exponencial
  dividida por uma potência de $x$ e multiplicada por um poliônimo.
  Ao calcularmos o limite $x\to 0^+$ o polinômio sempre produzirá um
  valor finito, enquanto
  \[
    \lim_{x\to 0^+}\frac{e^{\sfrac{-1}{x}}}{x^k} = 0 \qquad k = 1, 2, 3, \dots
  \]
  Concluímos que todas as derivadas de $f(x)$ são zero na origem.
  Assim todos os coeficientes da sua Série de Taylor são nulos e a soma da
  série é a função constante igual a zero, que, evidentemente, não é igual a $f$.
\end{example}
%------------------------------------------------------------------------------%

Uma função $f$ cuja Série de Taylor converge para $f$
é chamada de \emph{Função Analítica}\index{Função!analítica}. Nosso objetivo aqui
é poder determinar quais funções são analíticas.

O próximo exemplo mostra uma forma direta para demonstrarmos que a Série de Taylor
de uma função coincide com a função. Ainda nessa seção apresentaremos o Teorema de Taylor
que simplifica essa verificação.

%------------------------------------------------------------------------------%
\begin{example}
  Seja $f(x)=\ln(x)$ encontre a série de Taylor de $f$ centrada em $a=1$.
  Encontre o intervalo onde a série converge absolutamente e mostre
  que neste intervalo a série coincide com~$f$.
  \solution
  Note que
  \begin{align*}
    f(x)       & = \ln(x)            & f(1)       & =  0      \\[3mm]
    f^{(1)}(x) & = x^{-1}            & f^{(1)}(1) & =  1      \\[3mm]
    f^{(2)}(x) & = -x^{-2}           & f^{(2)}(1) & = -1      \\[3mm]
    f^{(3)}(x) & = 2x^{-3}           & f^{(3)}(1) & =  2      \\[3mm]
    f^{(4)}(x) & = -2\cdot 3\,x^{-4} & f^{(4)}(1) & = -6  \;.
  \end{align*}
  Podemos escrever, para todo $n\geq 1$, que
  \[
    f^{(n)}(x) = \frac{(-1)^{n-1}(n-1)!}{x^n}
  \]
  portanto
  \[
    f^{(n)}(1) = (-1)^{n-1}(n-1)! \qquad n\geq 1 \;.
  \]
  Então, a série de Taylor de $\ln(x)$ em $a=1$ é dada por
  \begin{align*}
    T(x)
    & = \sum_{n=0}^\infty \frac{f^{(n)}(1)}{n!}(x-1)^n \\[3mm]
    & = \frac{(-1)^0 f(1)}{0!}(x-1)^0 + \sum_{n=1}^\infty \frac{f^{(n)}(1)}{n!}(x-1)^n \\[3mm]
    & = \sum_{n=1}^\infty \frac{(-1)^{n-1}(n-1)!}{n!} (x-1)^n \\[3mm]
    & = \sum_{n=1}^\infty \frac{(-1)^{n-1}}{n} (x-1)^n  \;.
  \end{align*}
  Para encontrarmos o raio de convergência da série
  vamos aplicar o Teste da Raiz, ou seja, vamos calcular o limite
  \begin{align*}
    \sqrt[n]{\abs{a_n}}
    & = \sqrt[n]{\abs{\frac{(-1)^{n-1}}{n}(x-1)^n}} \\[3mm]
    & = \sqrt[n]{\frac{\abs{x-1}^n}{n}} \\[3mm]
    & = \frac{\abs{x-1}}{\sqrt[n]{n}} \\[3mm]
    & \rightarrow \abs{x-1}
  \end{align*}
  onde usamos que \(\sqrt[n]{n}\rightarrow 1\).
  Verificamos que o raio de convergência da série é $R=1$,
  portanto a série converge absolutamente para $x$ tal que $\abs{x-1}<1$,
  ou seja, para $x\in (0, 2)$.
  Poderíamos analisar os extremos desse intervalo,
  mas para identificar a série com a função $f$
  vamos precisar derivar a série termo a termo
  e esse procedimento só está garantido no intervalo aberto.

  Derivando \(T(x)\) temos
  \begin{align*}
    T'(x)
    & = \sum_{n=1}^\infty (-1)^{n-1}\frac{d}{dx}\frac{(x-1)^n}{n} \\[3mm]
    & = \sum_{n=1}^\infty (-1)^{n-1}(x-1)^{n-1} \\[3mm]
    & = \sum_{n=1}^\infty (1-x)^{n-1}    \;.
  \end{align*}
  Observe que essa última série é uma Série Geométrica com $\alpha=1$ e $r=1-x$.
  Como era de se esperar ela converge para $\abs{r}=\abs{1-x}<1$.
  Além disso, nesse intervalo sua soma é dada por
  \[
    T'(x) = \frac{\alpha}{1-r}
          = \frac{1}{1-(1-x)}
          = \frac{1}{\,x\,}   \;.
  \]
  Então, integrando $g'(t)$ sobre $[1, x]$ e usando que $g(1)=0$, obtemos
  \[
    T(x) = g(x) - g(1)
         = \int_1^x g'(t)\,dt
         = \int_1^x \frac{1}{\,t\,} dt
         = \ln(x)   \;.
  \]
  Dessa forma, fica provado que, no intervalo $(0,\,2)$, a função $\ln(x)$
  coincide com sua Série de Taylor centrada em $x=1$.

\end{example}
%------------------------------------------------------------------------------%

A Figura~\ref{fig:funcao_tayplor_logaritmo} mostra a função $\ln(x)$, em azul,
e seus polinômios de Taylor de ordens 1, 2, 3 e 6, o último desses polinômios
está desenhado em magenta. Observe que para $x$ entre zero e dois os polinômios
se aproximam da função logaritmo, porém a partir de $x=2$ eles se afastam
rapidamente.

\begin{figure}
  \begin{center}%
    \vspace{-0.2\baselineskip}%
    \scalebox{0.9}{% Title: gl2ps_renderer figure
% Creator: GL2PS 1.4.2, (C) 1999-2020 C. Geuzaine
% For: Octave
% CreationDate: Tue Sep 28 14:08:45 2021
\setlength{\unitlength}{1pt}
\begin{picture}(0,0)
\includegraphics[width=17cm,height=7cm]{fig_taylor_log}
\end{picture}%
\begin{picture}(481,198)(0,0)
\fontsize{10}{0}\selectfont\put(88.9306,137.666){\makebox(0,0)[b]{\textcolor[rgb]{0.15,0.15,0.15}{{0,5}}}}
\fontsize{10}{0}\selectfont\put(159.024,137.666){\makebox(0,0)[b]{\textcolor[rgb]{0.15,0.15,0.15}{{1}}}}
\fontsize{10}{0}\selectfont\put(229.117,137.666){\makebox(0,0)[b]{\textcolor[rgb]{0.15,0.15,0.15}{{1,5}}}}
\fontsize{10}{0}\selectfont\put(299.21,137.666){\makebox(0,0)[b]{\textcolor[rgb]{0.15,0.15,0.15}{{2}}}}
\fontsize{10}{0}\selectfont\put(369.303,137.666){\makebox(0,0)[b]{\textcolor[rgb]{0.15,0.15,0.15}{{2,5}}}}
\fontsize{10}{0}\selectfont\put(439.396,137.666){\makebox(0,0)[b]{\textcolor[rgb]{0.15,0.15,0.15}{{3}}}}
\fontsize{10}{0}\selectfont\put(13.8417,20.4678){\makebox(0,0)[r]{\textcolor[rgb]{0.15,0.15,0.15}{{-2}}}}
\fontsize{10}{0}\selectfont\put(13.8417,75.3265){\makebox(0,0)[r]{\textcolor[rgb]{0.15,0.15,0.15}{{-1}}}}
\fontsize{10}{0}\selectfont\put(13.8417,185.044){\makebox(0,0)[r]{\textcolor[rgb]{0.15,0.15,0.15}{{1}}}}
\fontsize{10}{0}\selectfont\put(457.62,119.214){\makebox(0,0)[l]{\textcolor[rgb]{0,0,0}{{$P_2$}}}}
\fontsize{10}{0}\selectfont\put(415.564,12.7876){\makebox(0,0)[l]{\textcolor[rgb]{0,0,0}{{$P_6$}}}}
\end{picture}
}%
    \vspace{-\baselineskip}%
  \end{center}
  \caption{Comparação da função logaritmo com seus polinômios de Taylor.}
  \label{fig:funcao_tayplor_logaritmo}
\end{figure}

O método usado no último exemplo funcionou bem para a função logaritmo,
mas buscar uma estratégia especifica para cada função pode ser
bastante complicado.
Felizmente, o Teorema de Taylor fornece uma forma relativamente
simples para verificar a convergência da série para uma
classe muito grande de funções.

A ideia do teorema é que podemos escrever uma função $f$ em
uma vizinhança de $x=a$ como sendo o seu polinômio de Taylor de
ordem $n$, $P_n(x)$, somado com o erro, $R_n(x)$, ou seja,
\[
  f(x) = P_n(x) + R_n(x)
       = \sum_{k=0}^n \frac{f^{(k)}(a)}{k!}(x-a)^k
       + R_n(x)    \;.
\]
Se formos capazes de encontrar os valores de $x$ para os quais
$R_n(x)$ tende a zero quando o grau do polinômio vai para infinito,
podemos concluir que $f$ coincide com sua Série de Taylor nestes valores.
Verificamos isso observando que
\begin{align*}
	f(x) & = \lim_{n\rightarrow \infty} \big( P_n(x) + R_n(x) \big)                    \\[3mm]
	     & = \lim_{n\rightarrow \infty} P_n(x) + \lim_{n\rightarrow \infty} R_n(x)     \\[3mm]
       & = \lim_{n\rightarrow \infty} P_n(x) + 0                                     \\[3mm]
       & = \lim_{n\rightarrow \infty} \sum_{k=0}^n \frac{f^{(k)}(a)}{k!}(x-a)^k      \\[3mm]
	     & = \sum_{k=0}^\infty \frac{f^{(k)}(a)}{k!}(x-a)^k    \;.
\end{align*}
A importância do Teorema de Taylor está exatamente no fato que ele fornece
uma fórmula para o resto $R_n(x)$.

%------------------------------------------------------------------------------%
\begin{teorema}{Teorema de Taylor}{teo:teorema-taylor}
  Seja%
  \index{Taylor!teorema}%
  \index{Teorema!Taylor}
  $f\colon I\rightarrow \R$ uma função que tem $n+1$ derivadas contínuas no
  intervalo aberto $I$, então para $a$, $x \in I$, podemos escrever que
  \[
    f(x) = P_n(x) + R_n(x)
         = \sum_{k=0}^n\frac{\;f^{(k)}(a)\;}{k!}(x-a)^k+ R_n(x)
  \]
  onde
  \[
    R_n(x)=\frac{\;f^{(n+1)}(c)\;}{(n+1)!}(x-a)^{n+1}
  \]
  para algum $c$ entre $a$ e $x$.
\end{teorema}
%------------------------------------------------------------------------------%
\begin{proof}
  Dado um $x$ fixo, que podemos considerar sem perda de generalidade menor do que $a$,
  definimos a constante $K$, de modo que
  \[
    f(x) = f(a)+f'(a)(x-a)
         + \cdots
         + \frac{f^{(n)}(a)}{n!}(x-a)^n
         + \frac{K}{(n+1)!}(x-a)^{n+1}    \;.
  \]
  Precisamos mostrar que existe $c\in(x, a)$ tal que
  \[
    K = f^{(n+1)}(c)    \;.
  \]
  Para isso definimos a função
  \begin{align*}
  	\phi(y) = f(x) - \bigg[\; f(y)   + f'(y)(x-y) + \cdots
  	                               & + \frac{f^{(n)}(y)}{n!}(x-y)^n        \\[2mm]
  	                               & +\frac{K}{(n+1)!}(x-y)^{n+1} \;\bigg]    \;.
  \end{align*}
  Por construção essa função é tal que
  \[
    \phi(a) = \phi(x) = 0    \;.
  \]
  Então, pelo Teorema de Rolle~\ref{teo:rolle},
  existe $c\in(x,\,a)$ tal que $\phi'(c)=0$.
  Calculando a derivada de $\phi(y)$ obtemos
  \begin{align*}
    \phi'(y) = -\bigg[ & \;\hpm f^{(1)}(y)                               \\[3mm]
                       & + f^{(2)}(y)(x-y) - f^{(1)}(y)                  \\[3mm]
                       & + \frac{f^{(3)}(y)}{2}(x-y)^2 - f^{(2)}(y)(x-y) \\[3mm]
                       & + \cdots                                        \\[3mm]
                       & + \frac{f^{(n+1)}(y)}{n!}   (x-y)^n
                         - \frac{f^{(n)  }(y)}{n!} n (x-y)^{n-1}         \\[3mm]
                       & -(n+1)\frac{K}{(n+1)!}(x-y)^{n}                 \;\bigg]    \;.
  \end{align*}
  Subtraindo os termos repetidos obtemos
  \[
    \phi'(y) = -\bigg[ \frac{f^{(n+1)}(y)}{n!}   (x-y)^n
                      -(n+1)\frac{K}{(n+1)!}(x-y)^{n}  \;\bigg]
  \]
  rearranjando podemos escrever
  \[
    \phi'(y) = \left(K - f^{(n+1)}(y)\right) \frac{(x-y)^{n}}{n!}    \;.
  \]
  Avaliando essa derivada em $c$ temos
  \[
    \phi'(c) = \left(K - f^{(n+1)}(c) \right) \frac{(x-c)^{n}}{n!} = 0    \;.
  \]
  Como $c\in(x, a)$, ou seja, $c\neq x$, segue que
  \[
    K = f^{(n+1)}(c)    \;.
  \]
\end{proof}
%------------------------------------------------------------------------------%

Para garantirmos que uma Série de Taylor converge para a sua função,
podemos mostrar que a diferença entre as somas parciais, $P_n(x)$,
e a função, $f$, tende à zero, para algum intervalo contendo $a$,
quando o número de termos vai para infinito, isso é,
\[
  \lim_{n\to \infty} \frac{\;f^{(n+1)}(c)\;}{(n+1)!}(x-a)^{n+1} = 0    \;.
\]
Porém, não existe uma forma sistemática para encontrar o valor correto para $c$,
o que normalmente fazemos é analisar todos os valores possíveis
e considerar o pior caso, o próximo resultado formaliza esse procedimento.

%------------------------------------------------------------------------------%
\begin{proposicao}{Desigualdade de Taylor}{prop:desigualdade-taylor}
  Se para \(\abs{x-a}\leq d\) vale a desigualdade
  \[
    \abs{f^{(n+1)}(x)} \leq M    \;.
  \]
  Então o erro da Série de Taylor de $f$ centrada em $a$
  satisfaz a \emph{Desigualdade de Taylor}%
  \index{Taylor!desigualdade}%
  \index{Desigualdade!Taylor}
  \[
    \abs{R_n(x)} \leq \frac{M}{\,(n+1)!\,}\abs{x-a}^{n+1}
  \]
  para todo $x$ tal que \(\abs{x-a}\leq d\).
\end{proposicao}
%------------------------------------------------------------------------------%

O próximo exemplo ilustra como podemos verificar quando a Série de Taylor
de uma função converge para a própria função.

%------------------------------------------------------------------------------%
\begin{example}
  Mostre que a Série de Taylor da função exponencial converge para a função
  exponencial.
  \solution
  Vimos no Exemplo~\ref{ex:serie-taylor-exponencial} que a Série de Taylor
  da função exponencial, $e^x$, centrada na origem, $a=0$, é convergente para
  todo $x\in\R$ e tem a forma
  \[
    T(x) = \sum_{n=0}^\infty \frac{x^n}{n!}
      = 1 + x + \frac{x^2}{2} + \frac{x^3}{3!} + \cdots + \frac{x^n}{n!}+\cdots
  \]
  Vamos demonstrar agora que a série converge para a função, isso é,
  \[
    e^x = \sum_{n=0}^\infty \frac{x^n}{n!}    \;.
  \]
  Para qualquer valor $x\in\R$ tomamos $d$ tal que
  \(
    \abs{x} \leq d
  \)
  temos então que para qualquer $n\in\N$
  \[
    \abs{f^{(n+1)}(x)} = e^x \leq e^d    \;.
  \]
  Podemos agora usar a Desigualdade de Taylor~\ref{prop:desigualdade-taylor}
  com $M=e^d$ obtendo
  \[
    \abs{R_n(x)} \leq \frac{e^d}{\,(n+1)!\,}\abs{x-a}^{n+1}    \;.
  \]
  Calculamos o limite encontramos
  \[
    \lim_{n\to \infty} \frac{e^d}{\,(n+1)!\,}\abs{x-a}^{n+1} =
    e^d \lim_{n\to \infty} \frac{\abs{x-a}^{n+1}}{\,(n+1)!\,} = 0    \;.
  \]
  Usando o Teorema do Confronto concluímos que
  \[
    \lim_{n\to \infty} R_n = 0
  \]
  e portanto a Série de Taylor converge para a função
  \[
    e^x = \sum_{n=0}^\infty \frac{x^n}{n!}    \;.
  \]
\end{example}
%------------------------------------------------------------------------------%

Temos agora os resultados necessários para demonstrar o cálculo da soma da
\emph{Série do Inverso do Fatorial}\index{Série!inverso do fatorial}
apresentada na Definição~\ref{def:serie-inverso-fatorial}.

%------------------------------------------------------------------------------%
\begin{example}[ex:soma-serie-inverso-fatorial]
  Mostre que a soma da Série do Inverso do Fatorial é
  \[
    \sum_{n=1}^\infty \frac{1}{\,n!\,} = e    \;.
  \]
  \solution
  Basta observarmos que a Série do Inverso do Fatorial coincide com a
  Série de Taylor da função exponencial em $x=1$
  \[
    \sum_{n=1}^\infty \frac{1  }{\,n!\,} =
    \sum_{n=1}^\infty \frac{1^n}{\,n!\,} = e^1 = e    \;.
  \]
\end{example}
%------------------------------------------------------------------------------%

A seguir apresentamos mais um exemplo de como podemos verificar que
a Série de Taylor converge para a sua função.

%------------------------------------------------------------------------------%
\begin{example}
  Encontre a série de Taylor de $\cos(x)$ centrada em $x=\sfrac{\pi}{2}$.
  Mostre que essa série converge para $\cos(x)$ para todo $x\in\R$.

  \solution
  Calculando as derivadas de $f(x) =\cos(x)$ obtemos
  \begin{align*}
    f^{(1)}(x) & =           -  \sin(x) \\
    f^{(2)}(x) & =           -  \cos(x) \\
    f^{(3)}(x) & = \hphantom{-} \sin(x) \;.
  \end{align*}
  Repetindo o processo percebemos que
  \begin{align*}
  	f^{(2n)}(x)   & =(-1)^{n}\hphantom{^{+1}} \cos(x)   \\
  	f^{(2n+1)}(x) & =(-1)^{n+1}               \sin(x)   \;.
  \end{align*}
  Calculando em $\sfrac{\pi}{2}$ temos
  \begin{align*}
     f^{(2n)}   & \left( \sfrac{\pi}{2} \right) = 0          \\
     f^{(2n+1)} & \left( \sfrac{\pi}{2} \right) = (-1)^{n+1} \;.
  \end{align*}
  Desse modo, a série de Taylor em $x=\sfrac{\pi}{2}$ é
  \[
    T(x) = \sum_{n=0}^\infty \frac{(-1)^{n+1}}{\,(2n+1)!\,}
    \left( x-\frac{\pi}{2} \right)^{2n+1}    \;.
  \]
  Pelo teorema de Taylor, sabemos que existe
  \(c \in \left(\sfrac{\pi}{2}, x \right)\), ou
  \(c \in \left(x, \sfrac{\pi}{2} \right)\), tal que
  \[
    \cos(x) = \sum_{n=0}^k \frac{(-1)^{n+1}}{\,(2n+1)!\,}
    \left(x-\frac{\pi}{2} \right)^{2n+1}
    + R_{2k+1}(x)
  \]
  onde
  \begin{align*}
   	R_{2k+1}(x) & =\frac{\, f^{(2k+2)}(c)     \,}{ (2k+2)! } \left( x - \frac{\pi}{2} \right)^{2k+2}  \\[3mm]
  	            & =\frac{\, (-1)^{n+1}\cos(c) \,}{ (2k+2)! } \left( x - \frac{\pi}{2} \right)^{2k+2}  \;.
  \end{align*}
  Como
  $\abs{\cos(c)} \leq 1$,
  segue que, para todo  $x\in \R$
  \[
    0 \;\leq\; \abs{R_{2k+1}(x)}
      \;\leq\; \frac{ \abs{x-\sfrac{\pi}{2}}^{2k+2} }{ (2k+2)! }
      \;\rightarrow\; 0   \;.
  \]
  Assim, $R_{2k+1}(x)\rightarrow 0$ para todo $x\in \R$. Isso implica que
  \[
    \cos(x) = \sum_{n=0}^\infty \frac{ (-1)^{n+1} }{\, (2n+1)! \,}
                                \left( x - \frac{\pi}{2} \right)^{2n+1}
    \qquad \forall\; x\in\R
  \]
\end{example}
%------------------------------------------------------------------------------%

\begin{figure}
  \begin{center}%
    \vspace{-0.2\baselineskip}%
    \scalebox{0.9}{% Title: gl2ps_renderer figure
% Creator: GL2PS 1.4.2, (C) 1999-2020 C. Geuzaine
% For: Octave
% CreationDate: Tue Sep 28 14:08:47 2021
\setlength{\unitlength}{1pt}
\begin{picture}(0,0)
\includegraphics[width=17cm,height=7cm]{fig_taylor_cos}
\end{picture}%
\begin{picture}(481,198)(0,0)
\fontsize{10}{0}\selectfont\put(19.2756,95.2752){\makebox(0,0)[t]{\textcolor[rgb]{0.15,0.15,0.15}{{$-3 \pi$}}}}
\fontsize{10}{0}\selectfont\put(91.5591,95.2752){\makebox(0,0)[t]{\textcolor[rgb]{0.15,0.15,0.15}{{$-2 \pi$}}}}
\fontsize{10}{0}\selectfont\put(163.843,95.2752){\makebox(0,0)[t]{\textcolor[rgb]{0.15,0.15,0.15}{{$-\pi$}}}}
\fontsize{10}{0}\selectfont\put(308.409,95.2752){\makebox(0,0)[t]{\textcolor[rgb]{0.15,0.15,0.15}{{$\pi$}}}}
\fontsize{10}{0}\selectfont\put(380.693,95.2752){\makebox(0,0)[t]{\textcolor[rgb]{0.15,0.15,0.15}{{$2 \pi$}}}}
\fontsize{10}{0}\selectfont\put(452.976,95.2752){\makebox(0,0)[t]{\textcolor[rgb]{0.15,0.15,0.15}{{$3 \pi$}}}}
\fontsize{10}{0}\selectfont\put(231.13,17.9742){\makebox(0,0)[r]{\textcolor[rgb]{0.15,0.15,0.15}{{-1}}}}
\fontsize{10}{0}\selectfont\put(231.13,187.538){\makebox(0,0)[r]{\textcolor[rgb]{0.15,0.15,0.15}{{1}}}}
\end{picture}
}%
    \vspace{-\baselineskip}%
  \end{center}
  \caption{Comparação da função cosseno com seus polinômios de Taylor.}
  \label{fig:funcao_tayplor_cosseno}
\end{figure}

A Figura~\ref{fig:funcao_tayplor_cosseno} mostra o gráfico do cosseno
em azul e os polinômios de Taylor
de graus 4, 8, 12 e 16, sendo que esse último destacado em magenta.

%------------------------------------------------------------------------------%

%------------------------------------------------------------------------------%

A Tabela~\ref{tab:series-maclaurin} lista algumas séries de Maclaurin
importantes e seus raios de convergência.

%------------------------------------------------------------------------------%
\begin{ntable}[\setlength{\extrarowheight}{5mm}]%
              {@{\extracolsep{-5pt}}%
               >{$\displaystyle\vphantom{\int_\int^\int}}w{c}{0.12\linewidth}<{$}%
               >{$\displaystyle\vphantom{\int_\int^\int}}w{l}{0.23\linewidth}<{$}%
               >{$\displaystyle\vphantom{\int_\int^\int}}w{l}{0.45\linewidth}<{$}%
               >{$\displaystyle\vphantom{\int_\int^\int}}w{l}{0.10\linewidth}<{$}}[1][5mm]
              {Séries de Maclaurin}
              {tab:series-maclaurin}
  \frac{1}{1-x}
  & = \sum_{n=0}^{\infty} x^n
  & = 1 + x + x^2 + x^3 + \cdots
  & R = 1
  \\ \hline
  e^x
  & = \sum_{n=0}^{\infty} \frac{x^n}{n!}
  & = 1 + \frac{x}{1}
    + \frac{x^2}{2!}
    + \frac{x^3}{3!}
    + \cdots
  & R = \infty
  \\ \hline
  \sin(x)
  & = \sum_{n=0}^{\infty} \frac{(-1)^n x^{2n+1}}{(2n+1)!}
  & = x - \frac{x^3}{3!}
    + \frac{x^5}{5!}
    - \frac{x^7}{7!}
    + \cdots
  & R = \infty
  \\ \hline
  \cos(x)
  & = \sum_{n=0}^{\infty} \frac{(-1)^nx^{2n}}{(2n)!}
  & = 1 - \frac{x^2}{2}
    + \frac{x^4}{4!}
    - \frac{x^6}{6!}
    + \cdots
  & R = \infty
  \\ \hline
  \arctan(x)
  & = \sum_{n=0}^{\infty} \frac{(-1)^n x^{2n+1}}{2n+1}
  & =	x - \frac{x^3}{3}
    + \frac{x^5}{5}
    - \frac{x^7}{7}
    + \cdots
  & R = 1
  \\ \hline
  \ln(x+1)
  & = \sum_{n=1}^{\infty} \frac{(-1)^{n-1}x^n}{n}
  & = x - \frac{x^2}{2}
    + \frac{x^3}{3}
    - \frac{x^4}{4}
    + \cdots
  & R = 1
  \\ \hline
  (1+x)^k
  & = \sum_{n=0}^{\infty} \binom{k}{n} x^n
  & = 1 + kx
    + \frac{k(k-1)}{2}x^2
    %	+ \frac{k(k-1)(k-2)}{3!}
    + \cdots \quad
  & R = 1
\end{ntable}
%------------------------------------------------------------------------------%

\FloatBarrier

%------------------------------------------------------------------------------%


%------------------------------------------------------------------------------%
